\section{方法}

\subsection{图与子集口径}

令人物图为 $G=(V,E)$：顶点 $V$ 是人物表中的实体，共 1231 个；边 $e\in E$ 是关系表中的一条记录，端点为关系两端的裸人物名，共 2324 条。边无向化（关系类别本身带方向语义，但本文只用量类别构成，方向不参与），自环与端点不在 $V$ 内的记录剔除。

子集取法同前，两条口径二选一，只取其一：

\textbf{口径 A（语义域命中）。} 给定词表 $W$，取人物表内在原始角色、身份、势力或别名字段中命中 $W$ 的人物构成 $S_A$；$|S_A|\ge 5$ 时用 A。

\textbf{口径 B（锚点自我网络）。} 给定锚点集 $A$，取所有至少一端属于 $A$ 的边构成 $E_B$，端点构成 $S_B$；$|S_A|<5$ 时启用。

本文该篇 $|S_A|<5$ 不满足口径 A 的下限，按回落规则改取口径 B（锚点自我网络）：以 4 位锚点各引其边，得 $|S_B|$=130 人、$E_B$=636 条；此后 $S$、$E(S)$ 均指该集合。

本文不采用「一跳邻域诱导子图」：预实验中它会退化成全图，把过半数人物与近九成的边一并卷进来，主题信号被淹没。

\subsection{类别构成与差异量}

对任一顶点集 $S$，令 $E(S)$ 为两端都在 $S$ 内的显式边，类别占比定义为

\begin{equation}\label{eq:p_cat}
p_c(S) = \frac{|\{e\in E(S): \mathrm{category}(e)=c\}|}{|E(S)|},
\end{equation}

基线取全图 $p_c(G)$。本文报告 $\Delta_c=p_c(S)-p_c(G)$，只看符号与量级；本子图显式边仅 636 条，显著性本身不可靠（第 6 节 L4），因此不给 $\Delta_c$ 配 p 值。

\subsection{热度与中心性}

对每位人物 $v$ 记 $h(v)$ 为其出场章节数、$d(v)$ 为其度数，取两者 Spearman 秩相关 $\rho$（并列秩取平均秩，p 值用正态近似）。这一步与主题无涉也照样算，是可复算量。

\subsection{可复算性}

本节全部量由 \texttt{.dump/}\_stats\_anthropology.json 提供，正文出现的每个小数都能在该文件中找到；作图脚本独立存于 figures/，重跑即可复现图 1--3。原书正文不外发，但命中计数由逐字匹配规则唯一确定，持有原书者可自行重算。
